% ============================================================
%  CHAPTER 6 — Vision-Language Reasoning
% ============================================================
\chapter{Vision-Language Reasoning}
\label{ch:vlm}

\section{Overview}

The reasoning backbone of the agent is \textbf{Google Gemini 2.0 Flash}, a multimodal
large language model accessed via the Google Generative Language REST API.
At each step of the control loop, the VLM receives the current camera image, a list
of YOLO detections, the robot's state, and the user's command, and responds with
a structured JSON object nominating which skill to execute next and why.

This chapter describes the VLM client implementation, the prompt design, the
structured output schema, and the retry and fallback mechanisms.

\section{Model Selection}

\begin{table}[H]
\centering
\caption{VLM selection criteria and outcome.}
\label{tab:vlm_select}
\begin{tabularx}{\linewidth}{lXX}
\toprule
\textbf{Criterion} & \textbf{Gemini 2.0 Flash} & \textbf{GPT-4o} \\
\midrule
Latency            & 0.4--1.0 s                & 0.8--2.0 s \\
Cost per image     & $\approx$\$0.005          & $\approx$\$0.015 \\
Python 3.6 compat. & Yes (REST only)            & Yes (REST only) \\
Context window     & 1M tokens                  & 128K tokens \\
Spatial reasoning  & Strong                     & Strong \\
JSON reliability   & High                       & High \\
\bottomrule
\end{tabularx}
\end{table}

Gemini 2.0 Flash was selected as the primary model due to its cost efficiency,
lower latency, and 1M-token context window, which allows embedding the full task
history and image without truncation.

\section{Implementation: VLMReasoner}

The \texttt{VLMReasoner} class encapsulates all API interactions.
It is initialised with a Google API key and communicates exclusively via the
\texttt{requests} library, making it compatible with Python 3.6.

\begin{lstlisting}[language=Python, caption={VLMReasoner initialisation.}]
class VLMReasoner:
    def __init__(self, api_key=None, model_name="gemini-2.0-flash"):
        self.api_key  = api_key or os.environ.get("EDENAI_API_KEY", "")
        self.api_url  = (
            "https://generativelanguage.googleapis.com/v1beta/models/"
            "{}:generateContent?key={}".format(model_name, self.api_key)
        )
        self.conversation_history = []   # rolling window of 20 turns
\end{lstlisting}

\section{Image Encoding}

Before sending to the API, each image undergoes three preprocessing steps:

\begin{enumerate}
  \item \textbf{Resize}: the longest dimension is clamped to 640 px, preserving the
        aspect ratio.
  \item \textbf{JPEG compression}: quality 85 — balances visual fidelity against
        token cost.
  \item \textbf{Base64 encoding}: the JPEG bytes are encoded to a UTF-8 string for
        inline embedding in the JSON payload.
\end{enumerate}

At 640$\times$480 JPEG quality 85, each image is approximately 40--60 KB, costing
roughly \$0.003--0.005 per API call.

\section{Prompt Design}
\label{sec:prompts}

\subsection{System Prompt}

The system prompt establishes the robot's identity, provides the current state,
lists available skills, and constrains the output format.
It is rebuilt at each query to embed the latest robot state.

\begin{mdframed}[style=promptbox]
\small
\begin{verbatim}
You are TIAGo, a mobile manipulator robot with an arm,
gripper, and omnidirectional base.

**Current State:**
- Gripper: {gripper_status}
- Base Location: {base_location}
- Detected Objects: {detected_objects}
- Recent Tasks: {last_3_tasks}

**Available Skills:**
grab_bottle() - Grasp any detected object (bottle, phone,
  cup, mug, etc.) in front of the robot using vision and
  arm control.
search_with_head(object_name='bottle') - Move head to scan
  surroundings for object.
handover_to_person(object='bottle') - Detect person with
  YOLO+depth, extend arm toward them, and hand over held
  object.
go_home() - Move arm to home/rest position.
wave() - Wave at person as a greeting gesture.
open_hand() - Open the robot gripper hand.
close_hand() - Close the robot gripper hand.

**Instructions:**
- Analyze the image to understand the visual scene.
- Choose ONE skill to execute next that makes progress
  toward the user's goal.
- Consider spatial relationships, object locations, and
  current robot state.
- If the object needed is NOT visible, use search_with_head
  FIRST before attempting to grab.
- Only choose grab_bottle if a graspable object is actually
  visible/detected.
- Explain your reasoning briefly (1-2 sentences).
- Respond ONLY with valid JSON in this exact format:
{
  "skill": "skill_name",
  "parameters": {},
  "rationale": "Why this skill makes sense"
}
Do not include any text before or after the JSON.
\end{verbatim}
\end{mdframed}

\subsection{User Prompt}

The user prompt appends the YOLO detections in a structured list format, with pixel
coordinates and confidence scores:

\begin{mdframed}[style=promptbox]
\small
\begin{verbatim}
User request: "Grab the bottle and give it to me"

Detected objects (YOLO):
- bottle @ bbox(142, 201, 85, 178) confidence=0.87
- person @ bbox(310, 120, 200, 380) confidence=0.93

What skill should I execute next?
\end{verbatim}
\end{mdframed}

\subsection{Multi-Step Continuation Prompt}

For the second and subsequent steps, the prompt is augmented with what has already
been done:

\begin{mdframed}[style=promptbox]
\small
\begin{verbatim}
Original task: "Grab the bottle and give it to me"
Completed so far: search_with_head, grab_bottle

If the original task is FULLY complete, respond with
skill='done'. Otherwise, what skill should execute next?
\end{verbatim}
\end{mdframed}

\subsection{Affordance Failure Prompt}

When a skill's affordance check fails, the failure reason is fed back to the VLM
so it can select a different skill:

\begin{mdframed}[style=promptbox]
\small
\begin{verbatim}
Original task: "Grab the bottle and give it to me"
Completed so far: nothing yet
Could NOT execute 'grab_bottle' because: No graspable
object detected in scene.
Choose a different skill to make progress (e.g.
search_with_head if object not visible).
\end{verbatim}
\end{mdframed}

\section{API Request Structure}

The Gemini REST API accepts a JSON payload with the image embedded inline.
System and user prompts are concatenated into a single text part:

\begin{lstlisting}[language=JSON, caption={Gemini API request payload structure.}]
{
  "contents": [{
    "parts": [
      {
        "text": "<SYSTEM PROMPT>\n\n<USER PROMPT>"
      },
      {
        "inline_data": {
          "mime_type": "image/jpeg",
          "data": "<BASE64_ENCODED_JPEG>"
        }
      }
    ]
  }]
}
\end{lstlisting}

\section{Response Parsing}

The API returns a JSON response from which the generated text is extracted:

\begin{lstlisting}[language=Python, caption={Response extraction and JSON parsing.}]
response_text = (result_data['candidates'][0]
                             ['content']['parts'][0]
                             ['text'].strip())

# Strip markdown code fences if present (Gemini sometimes wraps output)
if response_text.startswith("```"):
    response_text = response_text.split("```")[1]
    if response_text.startswith("json"):
        response_text = response_text[4:]
    response_text = response_text.strip()

result = json.loads(response_text)
# result = {"skill": "...", "parameters": {...}, "rationale": "..."}
\end{lstlisting}

\subsection{Output Schema}

\begin{lstlisting}[language=JSON, caption={VLM output JSON schema.}]
{
  "skill":      "grab_bottle",
  "parameters": {},
  "rationale":  "A bottle is visible at bbox(142,201,85,178)
                 with 87% confidence. Gripper is empty.
                 I will attempt to grasp it."
}
\end{lstlisting}

The \texttt{skill} field must match exactly one of the available skill names.
The \texttt{parameters} field carries optional skill-specific arguments
(\eg \texttt{\{"object\_name": "bottle"\}} for \texttt{search\_with\_head}).
The \texttt{rationale} is spoken aloud via TTS (truncated to 200 characters) and
logged for debugging.

\section{Retry Logic and Fallback}

The \texttt{query()} method retries up to 3 times on both JSON parse errors and
API communication errors:

\begin{lstlisting}[language=Python, caption={Retry logic.}]
for attempt in range(retry):      # retry = 3
    try:
        response = requests.post(self.api_url, json=payload, timeout=30)
        response.raise_for_status()
        result = json.loads(response_text)
        return result
    except json.JSONDecodeError:
        if attempt == retry - 1:
            return {'skill': 'go_home', 'parameters': {},
                    'rationale': 'VLM unavailable - safe default'}
        time.sleep(1)
    except Exception:
        if attempt == retry - 1:
            raise
        time.sleep(2)
\end{lstlisting}

On persistent failure after all retries, the safe fallback skill \texttt{go\_home}
is returned, ensuring the robot returns to a safe pose rather than taking no action.

\section{Post-Execution Verification}

After a skill executes, the VLM can optionally verify success by comparing before
and after images:

\begin{lstlisting}[language=Python, caption={Verification API call with two images.}]
payload = {
    "contents": [{
        "parts": [
            {"text": "BEFORE image:"},
            {"inline_data": {"mime_type": "image/jpeg", "data": pre_b64}},
            {"text": "AFTER image:"},
            {"inline_data": {"mime_type": "image/jpeg", "data": post_b64}},
            {"text": verification_prompt}
        ]
    }]
}
\end{lstlisting}

The verification prompt asks the VLM to respond with:
\begin{lstlisting}[language=JSON]
{
  "success": true,
  "observation": "The gripper is now closed around the bottle.",
  "issue": ""
}
\end{lstlisting}

\section{Conversation History}

The \texttt{VLMReasoner} maintains a rolling window of the last 20
\texttt{(user\_command, response)} pairs with timestamps.
This context is not currently injected into subsequent prompts (each query is
stateless at the API level) but is available for debugging and future extensions
where session memory is desired.

\section{Latency Analysis}

Gemini 2.0 Flash API latency depends on image size, prompt length, and network
conditions.
Table~\ref{tab:vlm_latency} shows typical values observed during testing.

\begin{table}[H]
\centering
\caption{Observed Gemini 2.0 Flash API latency.}
\label{tab:vlm_latency}
\begin{tabular}{ll}
\toprule
\textbf{Metric} & \textbf{Value} \\
\midrule
Minimum         & $\approx$ 0.4 s \\
Typical         & $\approx$ 0.7 s \\
Maximum (load)  & $\approx$ 1.5 s \\
Timeout         & 30 s \\
\bottomrule
\end{tabular}
\end{table}

This latency is acceptable for a non-real-time reasoning system where each VLM call
selects the next skill to execute (a decision that takes seconds to execute physically).
